% TOP_PROBLEM: {"id": 122, "title": "Covering Squares with Sumsets", "category_id": 1, "category": {"id": 1, "name": "number_theory", "display_name": "Number Theory", "description": "Properties of integers, prime numbers, Diophantine equations.", "slug": "number-theory", "order_index": 1, "created_at": "2026-07-31T15:26:25.670Z"}, "status": "open"}
% FIRST_POSTED: null
% ENTRY_KIND: statement_only
% Imported from: batch03.tex; UnsolvedMath ID 122
\documentclass[11pt]{article}
\usepackage[margin=25mm]{geometry}
\usepackage{fontspec}
\setmainfont{Latin Modern Roman}
\usepackage{amsmath,amssymb,mathtools}
\usepackage{unicode-math}
\setmathfont{Latin Modern Math}
\usepackage{xurl,hyperref}
\hypersetup{colorlinks=true,urlcolor=blue}
\setcounter{secnumdepth}{0}
\setlength{\parindent}{0pt}
\setlength{\parskip}{5pt}
\setlength{\emergencystretch}{3em}
\newcommand{\sref}[2]{\hyperlink{source:#1:#2}{[S#2]}}
\newcommand{\cataloguescope}[1]{\paragraph{Recorded target.}#1}
\begin{document}
% UnsolvedMath ID 122; source catalogue code GREEN-034
\setcounter{section}{121}
\section{Covering Squares with Sumsets}\label{122}
{\small\sffamily UnsolvedMath ID 122 · Number Theory}
\subsection{Definitions and mathematical statement}

\paragraph{Shared notation.}$\mathbb N=\{1,2,\ldots\}$, and $\mathbb Z,\mathbb Q,\mathbb R,\mathbb C$ denote the integers, rationals, reals and complex numbers. Any other conventions are specified locally.


For a finite set $A$ of real numbers, write $A+A=\{a+b:a,b\in A\}$; repeated summands are allowed. The first $n$ positive squares are $1^2,2^2,\ldots,n^2$. The asymptotic lower bound $n^{1-o(1)}$ means a lower bound $n^{1-\epsilon}$ for every fixed $\epsilon>0$, once $n$ is sufficiently large depending only on $\epsilon$.

The exponent question is equivalent if $A$ is required to consist of integers. Indeed, for arbitrary real $A$, the integer set $B=\{\lfloor a\rfloor,\lceil a\rceil:a\in A\}$ has $|B|\leq2|A|$, and every integer belonging to $A+A$ belongs to $B+B$. A constant factor of two does not affect the exponent $1-o(1)$.
\paragraph{Statement recorded in the supplied source.}
For every $\epsilon>0$, does there exist $n_0(\epsilon)$ such that for all integers $n\geq n_0(\epsilon)$ and all finite $A\subset\mathbb R$,
\[
 \{1^2,2^2,\ldots,n^2\}\subset A+A
 \quad\Longrightarrow\quad |A|\geq n^{1-\epsilon}?
\] \sref{122}{2}
\subsection{Short English statement}
Must a set whose pairwise sums contain every one of the first n positive squares have cardinality nearly n on the exponent scale?
\subsection{Sources}
\begin{description}
\item[\hypertarget{source:122:1}{S1}] ULAM.AI and the UnsolvedMath Contributors, \emph{UnsolvedMath: A Curated Collection of Open Mathematics Problems}, record 122 (GREEN-034). CC BY 4.0. \url{https://huggingface.co/datasets/ulamai/UnsolvedMath}.
\item[\hypertarget{source:122:2}{S2}] Ben Green, \emph{100 Open Problems}, maintained author collection, Problem 61 and its comments, accessed 7 October 2026. \url{https://people.maths.ox.ac.uk/greenbj/papers/open-problems.pdf}.
\end{description}
\label{end:122}
\end{document}
